\documentclass[../../../../main.tex]{subfiles}
\begin{document}

空間の3点$A$，$B$，$C$の組で，次の条件をみたすものを考える．
\begin{enumerate}
\item[\ding{"AC}]$A$，$B$，$C$は平面$x+y+z=1$上にある．
\item[\ding{"AD}]$A$の座標が$(l,m,n)$であれば，$B$，$C$の座標はそれぞれ$(m,n,l)$，$(n,l,m)$である．
\item[\ding{"AE}]ベクトル$\overrightarrow{OA}$，$\overrightarrow{OB}$は直交する(ただし，$O=(0,0,0)$)．
\end{enumerate}
　このとき，そのような$A$，$B$，$C$のとりかたに関せず次の三つが成り立つことを示せ。
\begin{enumerate}
  \item$A$，$B$，$C$は定円上にある．
  \item四面体$OABC$の体積は一定である．
  \item$BC$，$CA$，$AB$，$OA$，$OB$，$OC$の中点をそれぞれ$L$，$M$，$N$，$P$，$Q$，$R$とすれば
6点$L$，$M$，$N$，$P$，$Q$，$R$は定球面上にある．
\end{enumerate}

\end{document}
